\begin{lemma}[Claim 4.1(a2): controlled internal and tree-edge colors]
Let \(j\leq n\) be positive series indices.  Suppose \(\sigma_n\) is a
proper coloring leaving exactly \(e_0\) uncolored and is
\((T_n,D_n,\varphi_n)\)-stable.  If an exact edge identity \(f\) is
internal to \(T_j\),
\[
 \varphi_{j-1}(f)=c,
 \qquad c\in\overline\varphi_{j-1}(T_j)\cup D_{j-1},
\]
then
\[
                         \sigma_n(f)=\varphi_{j-1}(f).          \tag{1}
\]
In particular, (1) holds for every displayed edge identity of \(T_j\),
including its distinguished first edge \(e_0\).
\end{lemma}
\begin{proof}
Fix (f,c) as in the statement.  Every internal edge of (T_j) is
incident with (T_j), so the hypotheses of
\noderef{SameColorSeriesPreserved} apply at the initial index (j).  Use
its internal-edge conclusion over the interval through (n).  It gives
\[
                         \varphi_n(f)=\varphi_{j-1}(f)=c.       \tag{2}
\]
Iterating the two inclusions in
\noderef{HkuMissingStoredMonotone} along the ordered series gives
\[
 \overline\varphi_{j-1}(T_j)\cup D_{j-1}
 \subseteq \overline\varphi_n(T_n)\cup D_n.                    \tag{3}
\]
The same ordered initial-segment relations show that (f), already
incident with (T_j), is incident with (T_n).  By (2)--(3), its
\(\varphi_n\)-color is in the controlling set of the terminal stability
hypothesis.  The edge clause of \noderef{StableColoring} therefore gives
\(\sigma_n(f)=\varphi_n(f)\).  Combining this with (2) proves (1).

Now let (f) be a displayed edge of (T_j).  If it is the first edge,
\noderef{TreeSequence} identifies it with (e_0), and both colorings give
it value \(\mathrm{none}\) by \noderef{IsOneEdgeUncolored}.  Otherwise,
the first inclusion of \noderef{HkuTreeColorsControlled}, applied at
stage (j), says that its \(\varphi_{j-1}\)-color belongs to
\(\overline\varphi_{j-1}(T_j)\cup D_{j-1}\).  Its two endpoints are
displayed vertices by the attachment field of \noderef{TreeSequence}, so
the edge is internal to (T_j) according to \noderef{InternalEdges}.
The first part of the proof applies and yields (1).  This covers every
displayed edge identity and proves Claim 4.1(a2).
\end{proof}
